\bottomrule\end{tabularx}\endgroup\end{figure}

\clearpage
\section{Common conventions}
\hypertarget{tag-034E}{}\label{sec:II:conventions}
All spaces and coherent sheaves are complex analytic unless a polynomial
ring is specified. We use $B=\PP^1_t$, $(p_1,p_2,p_0)=(0,1,\infty)$,
the ordered lattice
$\Lambda=\Z\langle\widehat\gamma,\widehat u,\widehat w,\widehat\delta\rangle$,
and the marked periods
\[
\Pi_c=\begin{pmatrix}6\mu&\tau&1&0\\\beta^{\rm part}+c&\mu&0&1\end{pmatrix},
\qquad w=e^{-\pi i c},\qquad (\beta^{\rm part}+\tau)(p_0)=0.
\]
Base meridians are counterclockwise. The analytic root decks are clockwise,
$s_j(g_jz)=e^{-2\pi i/m_j}s_j(z)$, with $g_1g_2g_0=1$;
the resulting homology monodromies are $A_1,A_2,M_0$.
The primitive invariant vectors are
$\varepsilon=(1,2,-4,0)^t$ and
$\varepsilon'=(1,3,-3,0)^t$.
For a finite affine datum we write
$v_1=r_1\varepsilon+k_1\widehat\delta$ and
$v_2=r_2\varepsilon'+k_2\widehat\delta$; the coefficients
$r_1,r_2$ are also denoted $l_1,l_2$ in the peripheral presentation.
Define
\[
 D_{\rm part}=\operatorname{Im}\beta^{\rm part}
       -\frac{6(\operatorname{Im}\mu)^2}{\operatorname{Im}\tau},
 \qquad M=\max_{B\setminus\{p_0\}}D_{\rm part}
       =-\frac3\pi\log6.
\]
The admissible parameter inequality is $\operatorname{Im}c<-M$;
the maximum of $D_{\rm part}$ is attained only at $t=1$, as proved in \SeriesRef{fam:src:C.2.3}.
The standard multiplicities are $m_1=3,m_2=4$,
$f^*p_j=m_jS_j$, and $W=f^{-1}(p_0)$ is reduced, with normalization
$\eta:\mathrm dP_6\to W$.
For general finite affine twists we use
$p=12l_0-4l_1-3l_2$, retaining $\pi_1(X)=\Z/|p|$.
The complete marked transport and construction are those of Paper I.
An unqualified fibre is its full Cartier scheme; its reduction is always
specified explicitly. Homology and cohomology retain their displayed
coefficient groups, and complex manifolds carry their complex orientations.
